\documentclass[10pt,twocolumn]{article}
\usepackage[margin=0.75in]{geometry}
\usepackage{times}
\usepackage{amsmath,amssymb}
\usepackage{graphicx}
\usepackage{hyperref}
\usepackage{titlesec}
\usepackage{abstract}
\usepackage{enumitem}
\usepackage{fancyhdr}

\titleformat{\section}{\normalfont\large\bfseries}{\thesection}{1em}{}
\titleformat{\subsection}{\normalfont\bfseries}{\thesubsection}{1em}{}
\setlength{\columnsep}{0.3in}

\pagestyle{fancy}
\fancyhf{}
\renewcommand{\headrulewidth}{0pt}
\fancyfoot[C]{\thepage}

\title{\vspace{-0.5em}\LARGE\bfseries Systems Friction as a First-Class Design
Requirement for Outsiders in Civic Digital Services\\
\large Independent research note (portfolio writeup · not peer-reviewed)}

\author{
Jade Zhao\\
\small Informatics student, Luddy School (affiliation for contact only)\\
\small Indiana University Bloomington\\
\small \texttt{jlzhao@iu.edu}\\
\small\textit{Independent personal research ... not IU-endorsed or lab output}
}

\date{}

\begin{document}
\twocolumn[
  \begin{@twocolumnfalse}
    \maketitle
    \begin{abstract}
    \noindent
    Civic and public digital services often assume the user already has local
    context, language fluency, and unwritten process knowledge. For outsiders
    ... non-native speakers, new arrivals, and first-time users ... that
    assumption shows up as systems friction: stalled forms, PDF-only
    ``language support,'' blaming error messages, and channels that do not
    match outsider schedules. This independent research note locks a
    falsifiable claim that treating such friction as a first-class design
    requirement, rather than a post-launch accessibility fix, supports more
    accountable digital public services and public-interest AI adoption.
    Method: reflective pattern extraction from a Spring 2026 semester abroad
    in Madrid (Complutense / IU Overseas Study), documented as a public-safe
    pattern library, plus a reusable Outsider Impact Review one-pager for
    builders. Contribution at this stage is the claim, the patterns, the
    artifact template, and explicit limitations (single observer, one city,
    no controlled trial). No quantitative validation is claimed; a minimal
    kill test is specified for later work. This is a portfolio-style writeup,
    not a peer-reviewed publication, and it is not IU or lab research output.
    \end{abstract}
    \vspace{1em}
  \end{@twocolumnfalse}
]

\section{Introduction}
Public services increasingly run through forms, apps, and portals. When those
systems are designed for people who already know the local script ... which
office comes first, which document format is accepted, which language completes
the task ... newcomers experience the same interface as a wall. That wall is
easy to misread as personal failure. This note treats it as a design signal.

I spent Spring 2026 living and studying in Madrid through IU Overseas Study at
Universidad Complutense. Coursework supplied vocabulary for institutional
design; daily life supplied ground truth for how civic, university, commercial,
and transit systems meet someone who does not yet belong. Parallel clinic work
at ServeIT made the same failure modes visible in nonprofit handoffs: tools
built for the maintainer's mental model, not the visitor's first visit.

Existing accessibility practice (for example WCAG baselines used in clinic
work) catches many markup and interaction failures, but a happy-path or
checklist pass can still miss outsider-specific assumptions about local
context, sequencing, and language-as-access. Closely related discussions of
migrant administrative literacy, language barriers in e-government, and
linguistic fairness in AI document related harms; the gap this note targets is
narrower: a builder-facing requirement that outsider friction gate automation
and public digital changes \emph{before} launch, not only as a late
translation or accessibility patch.

This writeup's contributions, stated as claims:
\begin{itemize}[leftmargin=*,itemsep=2pt]
  \item Systems friction in civic digital services systematically disadvantages
        outsiders; the five patterns below are the working library from Madrid
        notes (suggestive, not causal proof).
  \item Treating that friction as a first-class design input ... via a reusable
        Outsider Impact Review ... is a practice bridge toward more accountable
        public digital services and public-interest AI adoption.
  \item The current method is reflective practice with an explicit kill test;
        no controlled experiment results are claimed here.
\end{itemize}
Section~2 situates prior themes; Section~3 describes method and design
principles; Section~4 reports observed patterns; Section~5 discusses limits;
Section~6 concludes.

\section{Related Work}
Three clusters frame the gap without claiming a complete literature survey
(deep reading remains in progress in the source repo's seed list).

\textbf{Accessibility and comprehensibility.} Clinic practice uses WCAG~2.1 as
a baseline for keyboard paths, contrast, and structure~\cite{wcag21}. Language
access and plain-language defaults sit adjacent to comprehensibility even when
checklists focus on markup. That overlap is necessary but not sufficient when
the primary task is only completable after hidden local prerequisites.

\textbf{Language, literacy, and civic access.} Public notes in this project
treat language choice as access control for task completion: Spanish-only
defaults, PDF-only ``English support,'' and mixed-language households that need
structural \texttt{lang} marking rather than a single page default. Adjacent
prior work on migrant administrative literacy and e-government language
barriers documents related burdens; this note does not claim to replace that
literature, only to reframe friction as a pre-automation design gate for
builders.

\textbf{Public-interest AI and accountable automation.} Serve-AI / ServeIT
clinic instincts (human review, honest handoffs, accessibility before
automation) rhyme with outsider friction as a ship/no-ship input. Policy and
fairness discussions of public-sector automated decision-making supply
accountability vocabulary; they rarely supply a portable one-pager students can
run before shipping clinic or civic automation. This note's artifact is aimed
at that operational gap.

\section{Approach / System}
Method now: reflective practice and pattern extraction from lived systems use
during one semester abroad, recorded as public-safe notes (no private
administrative detail). Exploration versus confirmation are kept separate in
the research log; confirmation experiments are planned, not run. The primary
artifact is an Outsider Impact Review one-pager: before shipping a civic flow
or adding automation, review assumed local context, unwritten sequencing,
language-as-access, hours and channels, recoverable errors, and automation
risk (whether AI would amplify those gaps).

Practice bridge only: ServeIT / Serve-AI clinic work is a rhyme, not a
relicensing of those materials. Serve-AI project materials are CC~BY-SA~4.0
with project credit; this madrid writeup and repo remain independent personal
notes under MIT.

\subsection{Design Principles}
\begin{itemize}[leftmargin=*,itemsep=2pt]
  \item Plain-language and completable language paths for the primary task
        (not PDF-only ``support'').
  \item Sequencing stated in the interface; no hidden prerequisite steps.
  \item Recoverable errors that name the fix and preserve state.
  \item Human escalation with case identity when the digital path fails.
  \item No assumed prior belonging: design for the first-time outsider, then
        add automation only if the access baseline works.
\end{itemize}

\section{Findings / Evaluation}
Patterns below appeared repeatedly in lived use of civic, university,
commercial, and transit/utility systems in Madrid. They \emph{suggest} design
failure modes; they do not prove causation beyond this observer's semester.
Recording method: public-safe reflective notes and essay outlines in the
\texttt{madrid-ai-ethics} repository.

\begin{enumerate}[leftmargin=*,itemsep=3pt]
  \item \textbf{Assumed local context.} Forms and apps expect address history,
        bank relationships, or national ID workflows a newcomer has never used.
  \item \textbf{Unwritten sequencing.} The flow works only if step~B happened
        before step~A, but the interface does not say so.
  \item \textbf{Language as access control.} Spanish-only defaults; ``English
        support'' that is a PDF rather than a completable interface; language
        choice functions like a login wall for task completion.
  \item \textbf{Hours and channels that do not match outsider life.} In-person
        only, narrow windows, little or no status tracking after submit ...
        especially hard for students and new arrivals.
  \item \textbf{Error messages that blame the user.} ``Invalid document''
        instead of naming the fix (for example a different scan format);
        recoverable error UX is missing when stakes are highest.
\end{enumerate}

Same failure modes show up in ServeIT clinic observation: tools built for the
maintainer's mental model, not the visitor's first visit. Evaluation bar for
later work (not yet run): apply Outsider Impact Review and a parallel
happy-path / WCAG-style pass to two public civic flows; if the outsider pass
finds nothing the happy-path pass misses, the ``first-class requirement'' half
of the claim is in serious trouble for this artifact. No baselines, sample
sizes, or percentage improvements are claimed here.

\section{Discussion}
Limitations are part of the claim package. Single observer: selection and
confirmation bias are real. One city and one semester (Madrid, Spring 2026):
other municipalities and seasons may differ. Reflective method: pattern
extraction from lived use, not lab measurement. Public-safe scope omits private
details that might strengthen or nuance a finding. Related-work deep reading is
still owned by the author; the seed list is not a finished survey. This note
does not claim legal or policy expertise, statistical generalizability, or that
automation always harms outsiders ... poorly scoped automation \emph{can}
amplify gaps; that is a design risk, not a theorem. It is not IRB, clinical, or
health-system research. Affiliation lines name where the author studies; they
do not mean Indiana University, ServeIT Clinic, or any lab endorses or owns
this content. Readers should not treat this PDF as a published conference paper
or as official IU research.

\section{Conclusion}
Outsiders hit civic digital systems where design assumes belonging. This
independent research note locks that friction as a first-class design concern,
documents five public-safe patterns from Madrid reflective practice, and offers
an Outsider Impact Review template for builders before automation. The next
empirical step is the dual-review kill test ... not invented metrics. Until
that runs, the contribution is the claim, the patterns, the artifact, and the
limits.

\section*{Acknowledgments}
Independent personal research notes. Drafting and structure were AI-assisted;
accountability for claims, limits, and publication decisions remains with Jade
Zhao. Practice bridge only to ServeIT Clinic and Serve-AI (PIT-UN); Serve-AI
project materials are CC~BY-SA~4.0 ... credit Serve-AI / ServeIT Clinic (IU
Luddy). This writeup is not Serve-AI licensed material and is not IU-endorsed
research.

\begin{thebibliography}{9}
\bibitem{wcag21}
W3C Web Accessibility Initiative,
\textit{Web Content Accessibility Guidelines (WCAG) 2.1},
W3C Recommendation, 5~June~2018.
\url{https://www.w3.org/TR/WCAG21/}.

\bibitem{serveit}
ServeIT Clinic, Luddy School of Informatics, Computing, and Engineering,
Indiana University.
\url{https://serveit.luddy.indiana.edu/}.

\bibitem{serveai}
Serve-AI project page (ServeIT Clinic / PIT-UN track).
\url{https://serveit.luddy.indiana.edu/serve-ai/index.html}.
\end{thebibliography}

\end{document}
